\section{Jin Architecture}


\subsection{Overview}

\begin{figure*}[t]
\centering
\begin{tikzpicture}[
    node distance=0.8cm,
    encoder/.style={rectangle, draw, minimum width=1.8cm, minimum height=1cm, align=center, fill=blue!10},
    joint/.style={rectangle, draw, minimum width=10cm, minimum height=1.2cm, align=center, fill=green!10},
    transformer/.style={rectangle, draw, minimum width=10cm, minimum height=1.5cm, align=center, fill=yellow!10},
    decoder/.style={rectangle, draw, minimum width=1.8cm, minimum height=1cm, align=center, fill=red!10},
    arrow/.style={->, thick}
]
% Inputs
\node[encoder] (text_enc) {Text\\Encoder};
\node[encoder, right=0.3cm of text_enc] (vision_enc) {Vision\\Encoder};
\node[encoder, right=0.3cm of vision_enc] (audio_enc) {Audio\\Encoder};
\node[encoder, right=0.3cm of audio_enc] (custom_enc) {Custom\\Encoder};

% Joint space
\node[joint, below=1cm of vision_enc] (joint) {Joint Embedding Space + Modality Tokens};

% Transformer
\node[transformer, below=0.8cm of joint] (transformer) {Diffusion Transformer with MoE};

% Decoders
\node[decoder, below=1cm of text_enc] (text_dec) at (text_enc |- transformer.south) {Text\\Output};
\node[decoder, below=1cm of vision_enc] (vision_dec) at (vision_enc |- transformer.south) {Vision\\Output};
\node[decoder, below=1cm of audio_enc] (audio_dec) at (audio_enc |- transformer.south) {Audio\\Output};
\node[decoder, below=1cm of custom_enc] (custom_dec) at (custom_enc |- transformer.south) {Custom\\Output};

% Arrows
\draw[arrow] (text_enc) -- (joint);
\draw[arrow] (vision_enc) -- (joint);
\draw[arrow] (audio_enc) -- (joint);
\draw[arrow] (custom_enc) -- (joint);
\draw[arrow] (joint) -- (transformer);
\draw[arrow] (transformer.south -| text_dec) -- (text_dec);
\draw[arrow] (transformer.south -| vision_dec) -- (vision_dec);
\draw[arrow] (transformer.south -| audio_dec) -- (audio_dec);
\draw[arrow] (transformer.south -| custom_dec) -- (custom_dec);

% Labels
\node[above=0.2cm of text_enc] {\small Tokens};
\node[above=0.2cm of vision_enc] {\small Patches};
\node[above=0.2cm of audio_enc] {\small Mel-Spec};
\node[above=0.2cm of custom_enc] {\small Features};
\end{tikzpicture}
\caption{Jin architecture: modality-specific encoders project into a joint embedding space, processed by a diffusion transformer with MoE, and decoded to target modalities.}
\label{fig:architecture}
\end{figure*}

Figure~\ref{fig:architecture} shows the Jin architecture. The system processes inputs through three stages:

\begin{enumerate}
    \item \textbf{Encoding}: Modality-specific encoders project inputs to the joint embedding space.
    \item \textbf{Processing}: A diffusion transformer with MoE processes the joint representation.
    \item \textbf{Decoding}: Modality-specific decoders generate outputs in target modalities.
\end{enumerate}

\subsection{Joint Embedding Space}

\begin{lstlisting}[language=Python]
class JointEmbeddingSpace(nn.Module):
    def __init__(self, config: JinConfig):
        self.encoders = nn.ModuleDict({
            "text": TextEncoder(config),
            "vision": VisionEncoder(config),
            "audio": AudioEncoder(config)
        })
        self.modality_embed = nn.Embedding(
            len(config.modalities), config.embedding_dim)
        self.fusion = nn.MultiheadAttention(
            config.embedding_dim, config.num_heads)

    def encode(self, inputs: Dict[str, Tensor]) -> Tensor:
        encoded = []
        for i, (mod, data) in enumerate(inputs.items()):
            embeds = self.encoders[mod].encode(data)
            mod_emb = self.modality_embed(
                torch.full((B, 1), i, device=device))
            embeds = embeds + mod_emb
            encoded.append(embeds)

        joint = torch.cat(encoded, dim=1)
        fused, _ = self.fusion(joint, joint, joint)
        return fused
\end{lstlisting}

\paragraph{Modality Tokens.} Each modality receives a learnable embedding added to all its tokens, enabling the transformer to distinguish modality origins without explicit conditioning.

\paragraph{Cross-Modal Fusion.} Self-attention over concatenated modality tokens enables cross-modal information flow from the first layer, unlike late-fusion approaches.

\subsection{Modality Encoders}

\paragraph{Text Encoder.} Standard token embedding with learned positional encodings:

\begin{lstlisting}[language=Python]
class TextEncoder(nn.Module):
    def encode(self, x: Tensor) -> Tensor:
        embeds = self.embed(x)  # (B, T, D)
        positions = self.pos_embed[:, :x.shape[1]]
        return embeds + positions
\end{lstlisting}

\paragraph{Vision Encoder.} Patch embedding following ViT~\cite{vit2020}:

\begin{lstlisting}[language=Python]
class VisionEncoder(nn.Module):
    def __init__(self, config, patch_size=14):
        self.proj = nn.Conv2d(3, config.embedding_dim,
                              kernel_size=patch_size,
                              stride=patch_size)

    def encode(self, x: Tensor) -> Tensor:
        patches = self.proj(x)  # (B, D, H/P, W/P)
        patches = rearrange(patches, 'b d h w -> b (h w) d')
        return patches + self.pos_embed[:, :patches.shape[1]]
\end{lstlisting}

\paragraph{Audio Encoder.} Projects mel-spectrogram frames:

\begin{lstlisting}[language=Python]
class AudioEncoder(nn.Module):
    def encode(self, x: Tensor) -> Tensor:
        # x: (B, T, 80) mel bins
        return self.proj(x) + self.pos_embed[:, :x.shape[1]]
\end{lstlisting}

\subsection{Diffusion Transformer with MoE}

\begin{lstlisting}[language=Python]
class DiffusionTransformerBlock(nn.Module):
    def __init__(self, config):
        self.norm1 = nn.LayerNorm(config.hidden_dim)
        self.attn = nn.MultiheadAttention(
            config.hidden_dim, config.num_heads)
        self.norm2 = nn.LayerNorm(config.hidden_dim)
        self.moe = UnifiedMoE(config)
        if config.use_diffusion:
            self.time_embed = nn.Sequential(
                nn.Linear(1, config.hidden_dim),
                nn.SiLU(),
                nn.Linear(config.hidden_dim,
                          config.hidden_dim))

    def forward(self, x, timestep=None):
        x = x + self.attn(self.norm1(x))[0]
        x = x + self.moe(self.norm2(x))
        if timestep is not None:
            x = x + self.time_embed(timestep).unsqueeze(1)
        return x
\end{lstlisting}

\paragraph{MoE Routing.} Each token is routed to top-$k$ experts:

\begin{lstlisting}[language=Python]
class UnifiedMoE(nn.Module):
    def forward(self, x: Tensor) -> Tensor:
        scores = self.router(x)  # (B, N, num_experts)
        weights, indices = torch.topk(
            F.softmax(scores, dim=-1),
            self.experts_per_token, dim=-1)

        output = torch.zeros_like(x)
        for i in range(self.experts_per_token):
            expert_idx = indices[:, :, i]
            expert_weight = weights[:, :, i].unsqueeze(-1)
            for e in range(self.num_experts):
                mask = (expert_idx == e)
                if mask.any():
                    expert_out = self.experts[e](x[mask])
                    output[mask] += expert_weight[mask] * expert_out
        return output
\end{lstlisting}

\subsection{z-JEPA: Hierarchical Zooming}

\begin{figure}[t]
\centering
\begin{tikzpicture}[
    node distance=0.6cm,
    level/.style={rectangle, draw, minimum width=2cm, minimum height=0.7cm, align=center},
    arrow/.style={->, thick}
]
\node[level, fill=blue!30] (l0) {Level 0\\224x224};
\node[level, below=of l0, fill=blue!25] (l1) {Level 1\\112x112};
\node[level, below=of l1, fill=blue!20] (l2) {Level 2\\56x56};
\node[level, below=of l2, fill=blue!15] (l3) {Level 3\\28x28};

\node[right=1.5cm of l0] (d0) {\small Global context};
\node[right=1.5cm of l1] (d1) {\small Object-level};
\node[right=1.5cm of l2] (d2) {\small Part-level};
\node[right=1.5cm of l3] (d3) {\small Fine detail};

\draw[arrow, bend right=20] (l0) to node[left] {\tiny predict} (l1);
\draw[arrow, bend right=20] (l1) to node[left] {\tiny predict} (l2);
\draw[arrow, bend right=20] (l2) to node[left] {\tiny predict} (l3);
\draw[arrow, bend left=20, dashed] (l3) to node[right] {\tiny feedback} (l2);
\draw[arrow, bend left=20, dashed] (l2) to node[right] {\tiny feedback} (l1);
\end{tikzpicture}
\caption{z-JEPA hierarchical levels with bidirectional prediction.}
\label{fig:zjepa}
\end{figure}

z-JEPA extends JEPA with hierarchical zoom levels (Figure~\ref{fig:zjepa}):

\begin{lstlisting}[language=Python]
@dataclass
class ZJEPAConfig:
    num_z_levels: int = 4
    zoom_factors: List[float] = [1.0, 0.5, 0.25, 0.125]
    predict_fine_from_coarse: bool = True
    predict_coarse_from_fine: bool = True
    bidirectional_prediction: bool = True
\end{lstlisting}

\paragraph{Cross-Scale Prediction.} The model learns to predict fine-grained features from coarse-grained context and vice versa:

\begin{equation}
    \mathcal{L}_{\text{cross}} = \sum_{i=0}^{L-1} \left( \|g_i(z_i) - z_{i+1}\|^2 + \|h_{i+1}(z_{i+1}) - z_i\|^2 \right)
\end{equation}

\paragraph{Adaptive Zoom.} An attention network identifies regions of interest for detailed processing:

\begin{lstlisting}[language=Python]
class AdaptiveZoomModule(nn.Module):
    def forward(self, x, global_features):
        roi_coords = self.roi_network(global_features)
        importance = self.importance_scorer(global_features)

        crops = [x]  # Level 0: full image
        for level in range(1, self.num_levels):
            crop = extract_foveal_region(
                x, roi_coords, level)
            crops.append(crop)
        return crops, importance
\end{lstlisting}

\subsection{Modality Extension}

New modalities can be added without retraining:

\begin{lstlisting}[language=Python]
class CustomModalityEncoder(nn.Module, ModalityEncoder):
    def __init__(self, input_dim, config):
        self.proj = nn.Linear(input_dim,
                              config.embedding_dim)

    def encode(self, x: Tensor) -> Tensor:
        return self.proj(x)

# Add at runtime
custom_encoder = CustomModalityEncoder(128, model.config)
model.add_modality("proprioception", custom_encoder)
\end{lstlisting}

The modality embedding table expands dynamically, preserving existing embeddings.

